\documentclass{subfiles}
\begin{document}
    The ElGamal schema is base on the hardness of the discrete logarithm, 
        i.e., the problem of finding the logarithm for \(a^{n} \mod n\),
        for some \(a, b \text{ and } n\).
    
    Algorithmically speaking we preoceed as follow: 
        we choose a cyclic group \(G\) generated by some integer \(\gamma\)
        such that \(\abs{G} = p\), with \(p\) a prime. 
        The group \(G^{*} = G \setminus \set{ 0 }\) is considered.
        Once the group is choosed, we create the public-keys as in \autoref{proc:ElGamal-keygen},
        and encode/decode the message using \autoref{proc:ElGamal-encode-decode}
        \subfile{../TikZ/Procedure 3 - ElGamal keygen}
        \subfile{../TikZ/Procedure 4 - ElGamal encode-decode.tex}

    The correctness of the decoding is trivial. Since the group is cyclic, 
        the product is commutative, meaning that \(\delta^{-a}\epsilon = \gamma^{-ak} m \gamma{ak} = m\).

\end{document}
